\documentclass[11pt]{article}
\usepackage[margin=1in]{geometry}
\usepackage{amsmath,amssymb,amsthm}
\usepackage[round]{natbib}
\usepackage{hyperref}
\usepackage{microtype}

\newtheorem{theorem}{Theorem}
\newtheorem{proposition}{Proposition}
\newtheorem{lemma}{Lemma}
\newtheorem{corollary}{Corollary}
\theoremstyle{remark}
\newtheorem{remark}{Remark}

\newcommand{\R}{\mathbb{R}}
\newcommand{\E}{\mathbb{E}}
\newcommand{\Prob}{\mathbb{P}}

\title{The Two Prices of Dependence:\\
a Vanishing Coverage Tax and a Fixed Sharpness Rent}
\author{Peter Cotton\\\texttt{peter.cotton@microprediction.com}}
\date{July 2026}

\begin{document}
\maketitle

\begin{abstract}
Temporal dependence charges split conformal prediction two different prices, and they scale
differently. The coverage price is a tax that vanishes with the calibration size: for a
stationary $\beta$-mixing score process the loss is at most
$\min_\tau\{\tau/(n+1)+2\beta(\tau)\}$ \citep{barberpananjady2026}. The sharpness price is a rent
that does not: the per-step log-score regret of the marginally calibrated predictive against the
past-conditional oracle equals the entropy-rate gap $\Delta H=H(S_0)-h$, the mutual information
between the present score and its infinite past, a constant independent of $n$. For Gaussian
score processes the rent exponentiates into width --- oracle intervals are narrower by the factor
$e^{-\Delta H}$ at matched coverage, which for an AR(1) score process with parameter $\phi$ is
$\sqrt{1-\phi^{2}}$. Dependence is therefore asymptotically free for the certificate and
permanently valuable for the forecast.
\end{abstract}

\section{Setting}\label{sec:setting}

Let $(S_t)_{t\in\mathbb Z}$ be a stationary ergodic sequence of conformity scores with continuous
marginal distribution $\bar F$ and marginal density $\bar f$. Write
$\mathcal P_t=\sigma(S_{t-1},S_{t-2},\dots)$ for the past, $F_t=\mathcal L(S_t\mid\mathcal P_t)$
for the past-conditional law with density $f(\cdot\mid\mathcal P_t)$, $H(S_0)$ for the
differential entropy of the marginal, and
$h=\lim_{m\to\infty}H(S_0\mid S_{-1},\dots,S_{-m})=H(S_0\mid\mathcal P_0)$ for the entropy rate
\citep{cover2006}. All entropies are assumed finite. Two predictive systems are compared at each
step:
\begin{itemize}
\item the \emph{marginal} system, which issues the stationary law $\bar F$ (this is what split
conformal calibrates: its quantile is an estimate of a quantile of $\bar F$, and the predictive
density it implies is the one residual shape $\bar f$, re-levelled);
\item the \emph{past-conditional oracle}, which issues $F_t$.
\end{itemize}
The comparison isolates what marginality itself forfeits; estimation error in either system is
not the subject.

\section{The coverage tax}\label{sec:tax}

Dependence costs the marginal system coverage, and the cost is quantified by
\citet{barberpananjady2026}: for a stationary $\beta$-mixing score sequence, split conformal with
$n$ calibration points and a memoryless pretrained score satisfies
\begin{equation}\label{eq:tax}
\Prob(\text{cover})\;\ge\;1-\alpha-\min_{\tau\ge 0}
\Bigl\{\tfrac{\tau}{n+1}+2\beta(\tau)\Bigr\},
\end{equation}
with a matching upper bound, and the bound is tight up to a constant factor over this class. For
fixed dependence the tax vanishes: for a process whose dependence dies at lag $t$
(e.g.\ MA($t$)), the loss is at most $(t+1)/(n+1)$, and for geometrically $\beta$-mixing
processes (e.g.\ stationary Gaussian AR(1)) it is $O(\log n/n)$. The earlier rate of
\citet{oliveira2024}, of order $\sqrt{\tau/n}$, is likewise vanishing. Nothing in this section is
new; it is quoted to fix one side of the comparison.

\section{The sharpness rent}\label{sec:rent}

\begin{proposition}[Rent identity]\label{prop:rent}
The per-step expected log-score regret of the marginal system against the past-conditional
oracle is
\[
\E\Bigl[\log\frac{f(S_0\mid\mathcal P_0)}{\bar f(S_0)}\Bigr]
\;=\;\E\,\mathrm{KL}\!\bigl(F_0\,\big\|\,\bar F\bigr)
\;=\;I\bigl(S_0;\mathcal P_0\bigr)
\;=\;H(S_0)-h\;=:\;\Delta H .
\]
\end{proposition}

\begin{proof}
The first equality is the tower property. The second is the standard identity
$I(X;Y)=\E_Y\,\mathrm{KL}(\mathcal L(X\mid Y)\,\|\,\mathcal L(X))$ applied with $X=S_0$,
$Y=\mathcal P_0$. The third is $I(S_0;\mathcal P_0)=H(S_0)-H(S_0\mid\mathcal P_0)$ and the
definition of the entropy rate of a stationary process.
\end{proof}

The rent is a property of the score process alone. It does not depend on the calibration size
$n$, and no re-levelling of the marginal shape can reduce it: conformalizing acts on $\bar f$ by
a monotone re-levelling, which leaves every likelihood ratio $f(\cdot\mid\mathcal P_0)/\bar
f(\cdot)$'s expectation unchanged \citep{cotton_marginally}. Proposition~\ref{prop:rent} is the
identity of that paper with the side information taken to be the past rather than a covariate;
the transfer requires stationarity and finite entropies, nothing else.

\begin{corollary}[Gaussian processes: rent as width]\label{cor:gauss}
Let $(S_t)$ be stationary Gaussian with marginal variance $\sigma^2$ and one-step innovation
variance $\sigma_\infty^2$ (the Kolmogorov--Szeg\H{o} prediction variance). Then
\[
\Delta H\;=\;\tfrac12\log\frac{\sigma^{2}}{\sigma_\infty^{2}},
\]
and at any matched coverage level the past-conditional oracle's central interval is narrower
than the marginal interval by the factor $\sigma_\infty/\sigma=e^{-\Delta H}$. For an AR(1)
process with parameter $\phi$ this factor is $\sqrt{1-\phi^{2}}$ and
$\Delta H=-\tfrac12\log(1-\phi^{2})$.
\end{corollary}

\begin{proof}
For Gaussian laws $H=\tfrac12\log(2\pi e\,\sigma^2)$ and $h=\tfrac12\log(2\pi e\,
\sigma_\infty^2)$; both central intervals are $z$-multiples of the relevant standard deviation at
matched level. For AR(1), $\sigma^2=\sigma_\infty^2/(1-\phi^2)$.
\end{proof}

\section{The asymmetry}\label{sec:asym}

Fix a stationary Gaussian AR(1) score process and let $n$ grow. The tax \eqref{eq:tax} is
$O(\log n/n)$; the rent is $-\tfrac12\log(1-\phi^2)$ at every $n$. Numerically:

\begin{center}
\begin{tabular}{c@{\qquad}c@{\qquad}c}
$\phi$ & rent $\Delta H$ (nats) & oracle width $/$ marginal width \\[2pt]
$0.3$ & $0.047$ & $0.954$ \\
$0.6$ & $0.223$ & $0.800$ \\
$0.9$ & $0.830$ & $0.436$ \\
\end{tabular}
\end{center}

At $\phi=0.9$ the marginal certificate costs essentially nothing to keep, a tax of order $\log n/n$, while the forecaster who conditions on the past issues intervals less than half as
wide at the same coverage, forever. The two prices answer different questions. The tax answers:
how much coverage does the certificate lose because the calibration data are dependent? The rent
answers: how much sharpness does the forecast lose because the calibration ignores that
dependence? The literature on conformal prediction for time series has concentrated on driving
the first price to zero; the second price is the one a forecaster pays at every step, and it is
not reduced by any amount of calibration data.

\section{Limits}\label{sec:limits}

Three restrictions delimit the claim. The comparison is between the marginal system and the
oracle, so it prices marginality itself; a feasible past-conditional model pays estimation costs
the oracle does not, and under strong dependence those costs can be material at small $n$. The
width statement of Corollary~\ref{cor:gauss} is Gaussian; for general processes the rent is the
log-score statement of Proposition~\ref{prop:rent}, and width comparisons depend on the shape of
the conditional laws. And the rent as stated prices the score process; when scores are residuals
from a fixed predictor, conditioning information in the original series that the predictor
already absorbs does not appear in $\Delta H$. Within these limits the accounting is exact, and
the direction of the asymmetry does not depend on the Gaussian case: the tax vanishes in $n$ and
the rent does not.

\bibliographystyle{plainnat}
\bibliography{../notes}

\end{document}
